Mauges Communauté, intercommunalité rurale et industrielle de 123~368 habitants située dans le Maine-et-Loire, produit chaque année avec ses 232~MW d'éoliennes et de panneaux solaires l'équivalent de 42~\% de l'électricité qu'elle consomme. Ce chiffre en volume masque pourtant une autre réalité, puisque la production et la consommation ne coïncident pas dans le temps. Ce rapport cherche jusqu'où le territoire peut développer le solaire et l'éolien sans perdre plus de 5~\% de sa production par écrêtement ni saturer son réseau, et avec quelle flexibilité. Aucune mesure horaire publique n'existant à cette échelle, nous avons reconstruit 43~824 heures de consommation et de production sur les années 2021 à 2025 à partir de données publiques de nature très diverse, électriques, météorologiques, démographiques et foncières. Chaque étape a été confrontée à des contrôles de cohérence. Un moteur de bilan horaire, mis à disposition des élus sous la forme d'un simulateur web, montre que doubler le solaire ne fait gagner que 8 points de couverture tout en faisant perdre 8,2~\% de la production. La batterie demandée par la fiche du projet supprime cet écrêtement mais coûte 186~M€ et ne se rembourse jamais. Nous recommandons plutôt d'augmenter l'éolien de 50~\% et le solaire de 25~\% en chauffant les ballons d'eau chaude à midi. La part de la consommation couverte heure par heure passe alors de 41,6 à 55,4~\%, sans stockage et sous la limite de 5~\% chaque année. Cette trajectoire n'est tenable qu'à condition d'orienter les nouveaux projets vers les postes sources qui disposent encore de capacité d'accueil, selon les capacités financières de raccordement disponibles de la région.

\vspace{0.3cm}
\noindent\textbf{Mots-clés~:} énergies renouvelables~; simulation horaire~; taux de couverture~; écrêtement~; flexibilité~; capacité d'accueil du réseau.
